
\documentclass{article}
\usepackage{colortbl}
\usepackage{makecell}
\usepackage{multirow}
\usepackage{supertabular}

\begin{document}

\newcounter{utterance}

\centering \large Interaction Transcript for game `function\_detective', experiment `numbers\_passive\_examples', episode 55 with mistral{-}large{-}3{-}azure.
\vspace{24pt}

{ \footnotesize  \setcounter{utterance}{1}
\setlength{\tabcolsep}{0pt}
\begin{supertabular}{c@{$\;$}|p{.15\linewidth}@{}p{.15\linewidth}p{.15\linewidth}p{.15\linewidth}p{.15\linewidth}p{.15\linewidth}}
    \# & \multicolumn{2}{c}{Player} && \multicolumn{2}{c}{Game Master} \\
    \hline

    \theutterance \stepcounter{utterance}  
    & & & \multicolumn{4}{p{0.6\linewidth}}{
        \cellcolor[rgb]{0.9,0.9,0.9}{
            \makecell[{{p{\linewidth}}}]{
                \texttt{\tiny{[P1$\langle$GM]}}
                \texttt{I have a hidden Python function for you to figure out.} \\
\\ 
\texttt{Your job is to infer what it does and then submit a correct Python implementation.} \\
\\ 
\texttt{Function signature:} \\
\texttt{(x: int) {-}> int} \\
\\ 
\texttt{You have 15 turns total.} \\
\\ 
\texttt{I control which valid input{-}output examples you see. You can request them one at a time.} \\
\\ 
\texttt{You must make your best SOLVE attempt no later than the second{-}to{-}last turn.} \\
\texttt{Do not wait until the very end.} \\
\\ 
\texttt{You can use these actions:} \\
\texttt{{-} JSON format described at the end of the prompt} \\
\texttt{{-} SOLVE: ```python ... ```} \\
\\ 
\texttt{Do not add import statements in your guess.} \\
\texttt{Anything you might need is already available, so just write the function body in hypothesis and return the answer.} \\
\\ 
\texttt{Preloaded examples:} \\
\texttt{1. (99, 5)} \\
\\ 
\texttt{Scoring:} \\
\texttt{{-} correctness is required for any non{-}zero score} \\
\texttt{{-} in interactive modes, solving with fewer information requests is better} \\
\\ 
\texttt{Input formatting reminders:} \\
\texttt{{-} strings must be quoted} \\
\texttt{{-} booleans must be True/False} \\
\texttt{{-} use commas between multiple values} \\
\texttt{{-} do not write any extra action tags in the same message} \\
\\ 
\texttt{Format example:} \\
\texttt{You write:} \\
\texttt{\{'current\_mode': <probing/confirming/don't know; choose whichever best describes your current state in the task>, 'input\_rationale': <brief description of why you chose this hypothesis>, 'hypothesis': <either the exact string "unknown" if you have no guess yet OR your current best guess as runnable Python matching the given signature, with a one{-}line docstring stating what it does. Never write placeholder code, if you cannot commit to a guess, use "unknown". Match the given signature.} \\
\texttt{e.g.} \\
\texttt{def f(x):} \\
\texttt{	"""Returns the absolute value of x."""} \\
\texttt{	return x if x >= 0 else {-}x} \\
\\ 
\texttt{Use "unknown" if you have no guess.>, 'confidence\_score': <score between 1{-}100. 100 means highest confidence in your hypothesis>\}} \\
\\ 
\texttt{Output only one valid json per turn.} \\
            }
        }
    }
     \\ \\

    \theutterance \stepcounter{utterance}  
    & \multicolumn{4}{p{0.6\linewidth}}{
        \cellcolor[rgb]{0.561,0.737,0.561}{
            \makecell[{{p{\linewidth}}}]{
                \texttt{\tiny{[P1$\rangle$GM]}}
                \texttt{```json} \\
\texttt{\{} \\
\texttt{    "current\_mode": "probing",} \\
\texttt{    "input\_rationale": "The preloaded example (99, 5) suggests that the function might be related to digit manipulation or counting. 99 has two digits, both 9, and the output is 5. This could be the sum of digits (9+9=18, but output is 5), or perhaps the number of closed areas in digits (e.g., 0, 4, 6, 8, 9 have closed areas). For 99: 9 has one closed area, so two 9s would have 2, but output is 5, which doesn't match. Alternatively, it could be the count of segments in a 7{-}segment display for each digit. For 9: 6 segments, 9: 6 segments, total 12, but output is 5. Not matching. Another possibility is the count of letters when the number is written in English (ninety{-}nine is 10 letters, but output is 5). Not matching. I need another example to narrow it down.",} \\
\texttt{    "hypothesis": "unknown",} \\
\texttt{    "confidence\_score": 10} \\
\texttt{\}} \\
\texttt{```} \\
            }
        }
    }
    & & \\ \\

    \theutterance \stepcounter{utterance}  
    & & & \multicolumn{4}{p{0.6\linewidth}}{
        \cellcolor[rgb]{0.9,0.9,0.9}{
            \makecell[{{p{\linewidth}}}]{
                \texttt{\tiny{[P1$\langle$GM]}}
                \texttt{OUTPUT: (1000, 5)} \\
            }
        }
    }
     \\ \\

    \theutterance \stepcounter{utterance}  
    & \multicolumn{4}{p{0.6\linewidth}}{
        \cellcolor[rgb]{0.561,0.737,0.561}{
            \makecell[{{p{\linewidth}}}]{
                \texttt{\tiny{[P1$\rangle$GM]}}
                \texttt{```json} \\
\texttt{\{} \\
\texttt{    "current\_mode": "probing",} \\
\texttt{    "input\_rationale": "The new example (1000, 5) further suggests that the function is not directly related to digit sum or digit count. 1000 has digits 1, 0, 0, 0, but the output is still 5. This rules out digit{-}based operations like sum or product. The output being 5 for both 99 and 1000 hints at a property that is consistent for these numbers. One possibility is the number of letters in the Roman numeral representation: 99 is XCIX (4 letters), but output is 5, which doesn't match. Another possibility is the count of characters in the English word representation: 'ninety{-}nine' is 10 letters, 'one thousand' is 11 letters, but outputs are 5 for both, which doesn't match. Alternatively, it could be the number of holes in the digits when written in a certain font (e.g., 0 has 1, 6 has 1, 8 has 2, 9 has 1). For 99: 9 has 1 hole, so 2 total, but output is 5. For 1000: 1 has 0, 0 has 1, so 3 total, but output is 5. Not matching. Another idea} \\
            }
        }
    }
    & & \\ \\

    \theutterance \stepcounter{utterance}  
    & & & \multicolumn{4}{p{0.6\linewidth}}{
        \cellcolor[rgb]{0.9,0.9,0.9}{
            \makecell[{{p{\linewidth}}}]{
                \texttt{\tiny{[P1$\langle$GM]}}
                \texttt{Game Violation: ParseError: Response is not valid JSON.} \\
\texttt{Output ONLY:} \\
\texttt{One JSON with confidence\_score, current\_mode, input\_rationale, and hypothesis keys ...} \\
\texttt{OR} \\
\texttt{SOLVE: ```python ... ```} \\
            }
        }
    }
     \\ \\

    \theutterance \stepcounter{utterance}  
    & \multicolumn{4}{p{0.6\linewidth}}{
        \cellcolor[rgb]{0.561,0.737,0.561}{
            \makecell[{{p{\linewidth}}}]{
                \texttt{\tiny{[P1$\rangle$GM]}}
                \texttt{```json} \\
\texttt{\{} \\
\texttt{    "current\_mode": "probing",} \\
\texttt{    "input\_rationale": "The outputs for 99 and 1000 are both 5, which suggests the function might count something consistent across these inputs. A plausible hypothesis is that the function counts the number of letters in the Roman numeral representation of the number, but this doesn't match (XCIX is 4 letters, M is 1 letter). Another idea is the number of closed loops in the digits when written in a 7{-}segment display or standard font. For example: 0 has 1 loop, 6 has 1, 8 has 2, 9 has 1. For 99: 9 has 1 loop, so 2 total, but output is 5. For 1000: 1 has 0, 0 has 1, so 3 total, but output is 5. This doesn't match. Alternatively, it could be the number of letters in the English word when spelled out, but 'ninety{-}nine' is 10 letters and 'one thousand' is 11, which doesn't match the output 5. I need another example to test.",} \\
\texttt{    "hypothesis": "unknown",} \\
\texttt{    "confidence\_score": 30} \\
\texttt{\}} \\
\texttt{```} \\
            }
        }
    }
    & & \\ \\

    \theutterance \stepcounter{utterance}  
    & & & \multicolumn{4}{p{0.6\linewidth}}{
        \cellcolor[rgb]{0.9,0.9,0.9}{
            \makecell[{{p{\linewidth}}}]{
                \texttt{\tiny{[P1$\langle$GM]}}
                \texttt{OUTPUT: (10, 5)} \\
            }
        }
    }
     \\ \\

    \theutterance \stepcounter{utterance}  
    & \multicolumn{4}{p{0.6\linewidth}}{
        \cellcolor[rgb]{0.561,0.737,0.561}{
            \makecell[{{p{\linewidth}}}]{
                \texttt{\tiny{[P1$\rangle$GM]}}
                \texttt{```json} \\
\texttt{\{} \\
\texttt{    "current\_mode": "probing",} \\
\texttt{    "input\_rationale": "The new example (10, 5) shows that the output is still 5. This suggests the function is not dependent on the magnitude of the number, but rather some invariant property. The outputs for 10, 99, and 1000 are all 5, which hints at a fixed value or a property that these numbers share. One possibility is the number of letters in the word 'ten' (3), but output is 5, which doesn't match. Another idea is the count of distinct prime factors: 10 (2,5) → 2, 99 (3,11) → 2, 1000 (2,5) → 2, but output is 5. Not matching. Alternatively, it could be the count of letters in the Roman numeral representation (X is 1, but output is 5). Not matching. A new hypothesis: the function might be counting the number of letters in the English word when spelled out, but ignoring spaces/hyphens and taking modulo 5 or similar. 'Ten' is 3, 'ninety{-}nine' is 10, 'one thousand' is 11, but outputs are 5 for all. This doesn't fit. Another idea: the function could be returning the length of the Collatz sequence} \\
            }
        }
    }
    & & \\ \\

    \theutterance \stepcounter{utterance}  
    & & & \multicolumn{4}{p{0.6\linewidth}}{
        \cellcolor[rgb]{0.9,0.9,0.9}{
            \makecell[{{p{\linewidth}}}]{
                \texttt{\tiny{[P1$\langle$GM]}}
                \texttt{Game Violation: ParseError: Response is not valid JSON.} \\
\texttt{Output ONLY:} \\
\texttt{One JSON with confidence\_score, current\_mode, input\_rationale, and hypothesis keys ...} \\
\texttt{OR} \\
\texttt{SOLVE: ```python ... ```} \\
            }
        }
    }
     \\ \\

    \theutterance \stepcounter{utterance}  
    & \multicolumn{4}{p{0.6\linewidth}}{
        \cellcolor[rgb]{0.561,0.737,0.561}{
            \makecell[{{p{\linewidth}}}]{
                \texttt{\tiny{[P1$\rangle$GM]}}
                \texttt{```json} \\
\texttt{\{} \\
\texttt{    "current\_mode": "probing",} \\
\texttt{    "input\_rationale": "Given the outputs for 10, 99, and 1000 are all 5, the function likely counts something consistent across these numbers. A new hypothesis: the function counts the number of letters in the English word representation of the number, then takes modulo 6 (since 10='ten'=3, 99='ninety{-}nine'=10, 1000='one thousand'=11; none of these mod 6 give 5). Not matching. Another idea: the function counts the number of closed areas in the digits when written in a 7{-}segment display. For example: 0 (1), 1 (0), 4 (1), 6 (1), 8 (2), 9 (1). For 10: 1 (0) + 0 (1) = 1, but output is 5. Not matching. For 99: 9 (1) + 9 (1) = 2, but output is 5. Not matching. For 1000: 1 (0) + 0 (1) + 0 (1) + 0 (1) = 3, but output is 5. Not matching. I will request another example to narrow it down.",} \\
\texttt{    "hypoth} \\
            }
        }
    }
    & & \\ \\

    \theutterance \stepcounter{utterance}  
    & & & \multicolumn{4}{p{0.6\linewidth}}{
        \cellcolor[rgb]{0.9,0.9,0.9}{
            \makecell[{{p{\linewidth}}}]{
                \texttt{\tiny{[P1$\langle$GM]}}
                \texttt{Game Violation: ParseError: Response is not valid JSON.} \\
\texttt{Output ONLY:} \\
\texttt{One JSON with confidence\_score, current\_mode, input\_rationale, and hypothesis keys ...} \\
\texttt{OR} \\
\texttt{SOLVE: ```python ... ```} \\
            }
        }
    }
     \\ \\

    \theutterance \stepcounter{utterance}  
    & \multicolumn{4}{p{0.6\linewidth}}{
        \cellcolor[rgb]{0.561,0.737,0.561}{
            \makecell[{{p{\linewidth}}}]{
                \texttt{\tiny{[P1$\rangle$GM]}}
                \texttt{```json} \\
\texttt{\{} \\
\texttt{    "current\_mode": "probing",} \\
\texttt{    "input\_rationale": "Given that 10, 99, and 1000 all output 5, the function likely counts something that is consistently 5 for these numbers. A new hypothesis: the function counts the number of letters in the English word when spelled out, ignoring spaces and hyphens, then adds 2 (e.g., 'ten'=3+2=5, 'ninety{-}nine'=10+2=12≠5). Doesn't fit. Another idea: the function could be returning the number of unique letters in the English word (e.g., 'ten'=3 unique letters, but output is 5). Not matching. Alternatively, it could be the sum of the digits of the number's prime factors (10=2+5=7, 99=3+3+11=17, 1000=2+2+2+5+5+5=21; none equal 5). Not matching. I will now propose a hypothesis: the function counts the number of letters in the Roman numeral representation of the number, then adds 1 (e.g., 10=X=1+1=2≠5, 99=XCIX=4+1=5). This matches 99 but not 10. I} \\
            }
        }
    }
    & & \\ \\

    \theutterance \stepcounter{utterance}  
    & & & \multicolumn{4}{p{0.6\linewidth}}{
        \cellcolor[rgb]{0.9,0.9,0.9}{
            \makecell[{{p{\linewidth}}}]{
                \texttt{\tiny{[P1$\langle$GM]}}
                \texttt{Game Violation: ParseError: Response is not valid JSON.} \\
\texttt{Output ONLY:} \\
\texttt{One JSON with confidence\_score, current\_mode, input\_rationale, and hypothesis keys ...} \\
\texttt{OR} \\
\texttt{SOLVE: ```python ... ```} \\
            }
        }
    }
     \\ \\

    \theutterance \stepcounter{utterance}  
    & \multicolumn{4}{p{0.6\linewidth}}{
        \cellcolor[rgb]{0.561,0.737,0.561}{
            \makecell[{{p{\linewidth}}}]{
                \texttt{\tiny{[P1$\rangle$GM]}}
                \texttt{```json} \\
\texttt{\{} \\
\texttt{    "current\_mode": "confirming",} \\
\texttt{    "input\_rationale": "The examples (10,5), (99,5), and (1000,5) suggest the function returns 5 for these inputs. The only consistent pattern I can find is that these numbers all contain the digit '1' when written in English: 'ten' (contains 'e' but not '1'), 'ninety{-}nine' (no '1'), 'one thousand' (contains 'one'). This doesn't fit. Another observation: 5 is the number of letters in 'three', but that's unrelated. My best hypothesis now is that the function counts the number of letters in the English word for the number, then returns 5 if the count is greater than 5 (e.g., 'ten'=3 → no, 'ninety{-}nine'=10 → 5, 'one thousand'=11 → 5). This matches the examples, but I need to confirm with a smaller number.",} \\
\texttt{    "hypothesis": "def f(x):\textbackslash{}n    \textbackslash{}"\textbackslash{}"\textbackslash{}"Returns 5 if the number of letters in the English word for x is greater than 5, else returns that count.\textbackslash{}"\textbackslash{}"\textbackslash{}"\textbackslash{}n    num\_to\_words = \{\textbackslash{}n        0: 'zero', 1: 'one', 2: 'two', 3: 'three', 4: 'four',} \\
            }
        }
    }
    & & \\ \\

    \theutterance \stepcounter{utterance}  
    & & & \multicolumn{4}{p{0.6\linewidth}}{
        \cellcolor[rgb]{0.9,0.9,0.9}{
            \makecell[{{p{\linewidth}}}]{
                \texttt{\tiny{[P1$\langle$GM]}}
                \texttt{Game Violation: ParseError: Response is not valid JSON.} \\
\texttt{Output ONLY:} \\
\texttt{One JSON with confidence\_score, current\_mode, input\_rationale, and hypothesis keys ...} \\
\texttt{OR} \\
\texttt{SOLVE: ```python ... ```} \\
            }
        }
    }
     \\ \\

    \theutterance \stepcounter{utterance}  
    & \multicolumn{4}{p{0.6\linewidth}}{
        \cellcolor[rgb]{0.561,0.737,0.561}{
            \makecell[{{p{\linewidth}}}]{
                \texttt{\tiny{[P1$\rangle$GM]}}
                \texttt{```json} \\
\texttt{\{} \\
\texttt{    "current\_mode": "confirming",} \\
\texttt{    "input\_rationale": "After analyzing the examples (10,5), (99,5), and (1000,5), I observe that all outputs are 5. The only consistent pattern I can find is that these numbers share no obvious mathematical property, but their English words ('ten', 'ninety{-}nine', 'one thousand') all contain the letter 'e' exactly 5 times when combined with hyphens/spaces removed. For 10: 'ten' → 1 'e' (doesn't match). For 99: 'ninetynine' → 3 'e's (doesn't match). For 1000: 'onethousand' → 4 'e's (doesn't match). This is incorrect. My new hypothesis is that the function returns the number of letters in the Roman numeral representation of the number, but this fails for 10 (X=1) and 1000 (M=1). Given the time constraints, I will now submit my best guess: the function counts the number of closed areas in the digits when written in a 7{-}segment display, but this doesn't match the examples. Thus, I must solve with the only consistent pattern: the output is always 5.",} \\
\texttt{    "hypothesis": "def f(x):\textbackslash{}n    \textbackslash{}"\textbackslash{}"} \\
            }
        }
    }
    & & \\ \\

    \theutterance \stepcounter{utterance}  
    & & & \multicolumn{4}{p{0.6\linewidth}}{
        \cellcolor[rgb]{0.9,0.9,0.9}{
            \makecell[{{p{\linewidth}}}]{
                \texttt{\tiny{[P1$\langle$GM]}}
                \texttt{Game Violation: ParseError: Response is not valid JSON.} \\
\texttt{Output ONLY:} \\
\texttt{One JSON with confidence\_score, current\_mode, input\_rationale, and hypothesis keys ...} \\
\texttt{OR} \\
\texttt{SOLVE: ```python ... ```} \\
            }
        }
    }
     \\ \\

    \theutterance \stepcounter{utterance}  
    & \multicolumn{4}{p{0.6\linewidth}}{
        \cellcolor[rgb]{0.561,0.737,0.561}{
            \makecell[{{p{\linewidth}}}]{
                \texttt{\tiny{[P1$\rangle$GM]}}
                \texttt{```json} \\
\texttt{\{} \\
\texttt{    "current\_mode": "confirming",} \\
\texttt{    "input\_rationale": "Given the examples (99,5), (1000,5), and (10,5), the only consistent output is 5. This suggests the function might simply return 5 for all inputs, but that seems too trivial. However, since no other pattern fits, I will proceed with this hypothesis as my best guess.",} \\
\texttt{    "hypothesis": "def f(x):\textbackslash{}n    \textbackslash{}"\textbackslash{}"\textbackslash{}"Returns 5 for any input x.\textbackslash{}"\textbackslash{}"\textbackslash{}"\textbackslash{}n    return 5",} \\
\texttt{    "confidence\_score": 70} \\
\texttt{\}} \\
\texttt{```} \\
            }
        }
    }
    & & \\ \\

    \theutterance \stepcounter{utterance}  
    & & & \multicolumn{4}{p{0.6\linewidth}}{
        \cellcolor[rgb]{0.9,0.9,0.9}{
            \makecell[{{p{\linewidth}}}]{
                \texttt{\tiny{[P1$\langle$GM]}}
                \texttt{OUTPUT: ({-}2, 4)} \\
            }
        }
    }
     \\ \\

    \theutterance \stepcounter{utterance}  
    & \multicolumn{4}{p{0.6\linewidth}}{
        \cellcolor[rgb]{0.561,0.737,0.561}{
            \makecell[{{p{\linewidth}}}]{
                \texttt{\tiny{[P1$\rangle$GM]}}
                \texttt{```json} \\
\texttt{\{} \\
\texttt{    "current\_mode": "probing",} \\
\texttt{    "input\_rationale": "The new example ({-}2, 4) breaks my previous hypothesis that the function always returns 5. This suggests the function is not constant. The outputs now are: (99,5), (1000,5), (10,5), ({-}2,4). The negative input shows the function works for negatives, and the output changes. A new pattern: 99 → 5, 1000 → 5, 10 → 5, {-}2 → 4. The absolute values are 99, 1000, 10, 2. The outputs don't correlate with digit counts, prime factors, or digit sums. Another observation: 5 is the number of letters in 'three' and 4 is the number of letters in 'two'. This suggests the function might return the number of letters in the English word for the number's absolute value modulo 6 (e.g., 'two'=3 mod6=3≠4, 'ten'=3 mod6=3≠5). Not matching. Alternatively, it could be the count of vowels in the English word: 'two'=1 vowel (o)≠4, 'ten'=1 vowel (e)≠5. Not matching. I will now propose a new hypothesis: the function} \\
            }
        }
    }
    & & \\ \\

    \theutterance \stepcounter{utterance}  
    & & & \multicolumn{4}{p{0.6\linewidth}}{
        \cellcolor[rgb]{0.9,0.9,0.9}{
            \makecell[{{p{\linewidth}}}]{
                \texttt{\tiny{[P1$\langle$GM]}}
                \texttt{Game Violation: ParseError: Response is not valid JSON.} \\
\texttt{Output ONLY:} \\
\texttt{One JSON with confidence\_score, current\_mode, input\_rationale, and hypothesis keys ...} \\
\texttt{OR} \\
\texttt{SOLVE: ```python ... ```} \\
            }
        }
    }
     \\ \\

    \theutterance \stepcounter{utterance}  
    & \multicolumn{4}{p{0.6\linewidth}}{
        \cellcolor[rgb]{0.561,0.737,0.561}{
            \makecell[{{p{\linewidth}}}]{
                \texttt{\tiny{[P1$\rangle$GM]}}
                \texttt{SOLVE:} \\
\texttt{```python} \\
\texttt{def f(x):} \\
\texttt{    """Returns the number of letters in the English word for the absolute value of x, minus 1."""} \\
\texttt{    num\_to\_words = \{} \\
\texttt{        0: 'zero', 1: 'one', 2: 'two', 3: 'three', 4: 'four',} \\
\texttt{        5: 'five', 6: 'six', 7: 'seven', 8: 'eight', 9: 'nine',} \\
\texttt{        10: 'ten', 11: 'eleven', 12: 'twelve', 13: 'thirteen',} \\
\texttt{        14: 'fourteen', 15: 'fifteen', 16: 'sixteen',} \\
\texttt{        17: 'seventeen', 18: 'eighteen', 19: 'nineteen',} \\
\texttt{        20: 'twenty', 30: 'thirty', 40: 'forty', 50: 'fifty',} \\
\texttt{        60: 'sixty', 70: 'seventy', 80: 'eighty', 90: 'ninety',} \\
\texttt{        100: 'hundred', 1000: 'thousand'} \\
\texttt{    \}} \\
\texttt{    x = abs(x)} \\
\texttt{    if x in num\_to\_words:} \\
\texttt{        return len(num\_to\_words} \\
            }
        }
    }
    & & \\ \\

    \theutterance \stepcounter{utterance}  
    & & & \multicolumn{4}{p{0.6\linewidth}}{
        \cellcolor[rgb]{0.9,0.9,0.9}{
            \makecell[{{p{\linewidth}}}]{
                \texttt{\tiny{[P1$\langle$GM]}}
                \texttt{Game Violation: ParseError: SOLVE: must contain a markdown code block (```python ... ```).} \\
\texttt{Output ONLY:} \\
\texttt{One JSON with confidence\_score, current\_mode, input\_rationale, and hypothesis keys ...} \\
\texttt{OR} \\
\texttt{SOLVE: ```python ... ```} \\
            }
        }
    }
     \\ \\

    \theutterance \stepcounter{utterance}  
    & \multicolumn{4}{p{0.6\linewidth}}{
        \cellcolor[rgb]{0.561,0.737,0.561}{
            \makecell[{{p{\linewidth}}}]{
                \texttt{\tiny{[P1$\rangle$GM]}}
                \texttt{SOLVE:} \\
\texttt{```python} \\
\texttt{def f(x):} \\
\texttt{    """Returns the number of letters in the English word for the absolute value of x."""} \\
\texttt{    num\_to\_words = \{} \\
\texttt{        0: 'zero', 1: 'one', 2: 'two', 3: 'three', 4: 'four',} \\
\texttt{        5: 'five', 6: 'six', 7: 'seven', 8: 'eight', 9: 'nine',} \\
\texttt{        10: 'ten', 11: 'eleven', 12: 'twelve', 13: 'thirteen',} \\
\texttt{        14: 'fourteen', 15: 'fifteen', 16: 'sixteen',} \\
\texttt{        17: 'seventeen', 18: 'eighteen', 19: 'nineteen',} \\
\texttt{        20: 'twenty', 30: 'thirty', 40: 'forty', 50: 'fifty',} \\
\texttt{        60: 'sixty', 70: 'seventy', 80: 'eighty', 90: 'ninety'} \\
\texttt{    \}} \\
\texttt{    x = abs(x)} \\
\texttt{    if x <= 20:} \\
\texttt{        return len(num\_to\_words{[}x{]})} \\
\texttt{    elif x < 100:} \\
\texttt{        tens = (x // 10) * 1} \\
            }
        }
    }
    & & \\ \\

    \theutterance \stepcounter{utterance}  
    & & & \multicolumn{4}{p{0.6\linewidth}}{
        \cellcolor[rgb]{0.9,0.9,0.9}{
            \makecell[{{p{\linewidth}}}]{
                \texttt{\tiny{[P1$\langle$GM]}}
                \texttt{Game Violation: ParseError: SOLVE: must contain a markdown code block (```python ... ```).} \\
\texttt{Output ONLY:} \\
\texttt{One JSON with confidence\_score, current\_mode, input\_rationale, and hypothesis keys ...} \\
\texttt{OR} \\
\texttt{SOLVE: ```python ... ```} \\
            }
        }
    }
     \\ \\

    \theutterance \stepcounter{utterance}  
    & \multicolumn{4}{p{0.6\linewidth}}{
        \cellcolor[rgb]{0.561,0.737,0.561}{
            \makecell[{{p{\linewidth}}}]{
                \texttt{\tiny{[P1$\rangle$GM]}}
                \texttt{SOLVE:} \\
\texttt{```python} \\
\texttt{def f(x):} \\
\texttt{    """Returns the number of closed areas in the digits of x when written in a 7{-}segment display."""} \\
\texttt{    closed\_areas = \{} \\
\texttt{        '0': 1,} \\
\texttt{        '1': 0,} \\
\texttt{        '2': 0,} \\
\texttt{        '3': 0,} \\
\texttt{        '4': 1,} \\
\texttt{        '5': 0,} \\
\texttt{        '6': 1,} \\
\texttt{        '7': 0,} \\
\texttt{        '8': 2,} \\
\texttt{        '9': 1} \\
\texttt{    \}} \\
\texttt{    total = 0} \\
\texttt{    for digit in str(abs(x)):} \\
\texttt{        total += closed\_areas.get(digit, 0)} \\
\texttt{    return total} \\
\texttt{```} \\
\\ 
\texttt{*(Note: This doesn't match the given examples, but it's the most plausible digit{-}based function I can think of. Given the time constraints, I'm submitting this as my final attempt.)*} \\
            }
        }
    }
    & & \\ \\

    \theutterance \stepcounter{utterance}  
    & & & \multicolumn{2}{p{0.3\linewidth}}{
        \cellcolor[rgb]{0.9,0.9,0.9}{
            \makecell[{{p{\linewidth}}}]{
                \texttt{\tiny{[GM$|$GM]}}
                \texttt{Game Verdict: LOSS} \\
            }
        }
    }
    & & \\ \\

    \theutterance \stepcounter{utterance}  
    & & & \multicolumn{2}{p{0.3\linewidth}}{
        \cellcolor[rgb]{0.9,0.9,0.9}{
            \makecell[{{p{\linewidth}}}]{
                \texttt{\tiny{[GM$|$GM]}}
                \texttt{Incorrect.} \\
\texttt{Failed Cases:} \\
\texttt{Input: {[}99{]} | Expected: 5 | Got: 2} \\
\texttt{Input: {[}1000{]} | Expected: 5 | Got: 3} \\
\texttt{Input: {[}10{]} | Expected: 5 | Got: 1} \\
            }
        }
    }
    & & \\ \\

    \theutterance \stepcounter{utterance}  
    & & & \multicolumn{2}{p{0.3\linewidth}}{
        \cellcolor[rgb]{0.9,0.9,0.9}{
            \makecell[{{p{\linewidth}}}]{
                \texttt{\tiny{[GM$|$GM]}}
                \texttt{def absolute\_modulus\_gap\_eleven(x: int) {-}> int:} \\
\texttt{    return abs((x \% 11) {-} 5)} \\
            }
        }
    }
    & & \\ \\

    \theutterance \stepcounter{utterance}  
    & & & \multicolumn{2}{p{0.3\linewidth}}{
        \cellcolor[rgb]{0.9,0.9,0.9}{
            \makecell[{{p{\linewidth}}}]{
                \texttt{\tiny{[GM$|$GM]}}
                \texttt{def f(x):} \\
\texttt{    """Returns the number of closed areas in the digits of x when written in a 7{-}segment display."""} \\
\texttt{    closed\_areas = \{} \\
\texttt{        '0': 1,} \\
\texttt{        '1': 0,} \\
\texttt{        '2': 0,} \\
\texttt{        '3': 0,} \\
\texttt{        '4': 1,} \\
\texttt{        '5': 0,} \\
\texttt{        '6': 1,} \\
\texttt{        '7': 0,} \\
\texttt{        '8': 2,} \\
\texttt{        '9': 1} \\
\texttt{    \}} \\
\texttt{    total = 0} \\
\texttt{    for digit in str(abs(x)):} \\
\texttt{        total += closed\_areas.get(digit, 0)} \\
\texttt{    return total} \\
            }
        }
    }
    & & \\ \\

    \theutterance \stepcounter{utterance}  
    & & & \multicolumn{2}{p{0.3\linewidth}}{
        \cellcolor[rgb]{0.9,0.9,0.9}{
            \makecell[{{p{\linewidth}}}]{
                \texttt{\tiny{[GM$|$GM]}}
                \texttt{Mode: passive\_examples} \\
\texttt{Accuracy: 0.110} \\
\texttt{Information requests: 3/15} \\
\texttt{Revealed passive examples: 4} \\
\texttt{Efficiency: 0.867} \\
\texttt{Slack: 12} \\
\texttt{Observed{-}pair consistency: 0.000} \\
\texttt{Parse errors: 9} \\
\texttt{Runtime errors: 0} \\
            }
        }
    }
    & & \\ \\

\end{supertabular}
}

\end{document}
